%% fema_peltzman — paper/draft.tex
%% Compile: pdflatex draft.tex (or: pandoc --pdf-engine=tectonic draft.tex -o draft.pdf)
%%
%% PLACEHOLDER SECTIONS: all tables, figures, and numeric results are imported
%% from the analysis pipeline outputs (tables/*.tex, figures/*.png).
%% Do NOT hardcode any quantitative values in this file.

\documentclass[12pt]{article}

%% --- Packages ---
\usepackage[margin=1in]{geometry}
\usepackage{setspace}
\doublespacing
\usepackage{amsmath,amssymb}
\usepackage{graphicx}
\usepackage{booktabs}
\usepackage{natbib}
\bibpunct{(}{)}{;}{a}{,}{,}
\usepackage{hyperref}
\hypersetup{colorlinks=true, linkcolor=blue, citecolor=blue, urlcolor=blue}
\usepackage{caption}
\usepackage{subcaption}
\usepackage{threeparttable}
\usepackage{rotating}
\usepackage{lmodern}
\usepackage[T1]{fontenc}

%% Figure/table path
\graphicspath{{../figures/}}

%% Convenience macros
\newcommand{\IA}{\textrm{IA}}
\newcommand{\NFIP}{\textrm{NFIP}}
\newcommand{\RD}{\textrm{RD}}
\newcommand{\IV}{\textrm{IV}}
\newcommand{\OLS}{\textrm{OLS}}
\newcommand{\SDID}{\textrm{SDID}}
\newcommand{\CCT}{\textrm{CCT}}

% ============================================================
\begin{document}
% ============================================================

%% Title page
\title{Does Disaster Aid Crowd Out Private Mitigation?\\
       Evidence from FEMA Individual Assistance Thresholds}
\author{Ryan Vaughn \and [Co-author TBD]}
\date{\today}
\maketitle

\begin{abstract}
We test whether exposure to FEMA Individual Assistance (\IA{}) alters household
pre-disaster mitigation investment and post-disaster claiming behavior—the
``Peltzman hypothesis'' applied to federal disaster aid.
Using the per-capita damage threshold that governs \IA{} eligibility as a
source of quasi-random variation, we estimate two complementary designs.
\emph{Mechanism A}: a historical instrumental variables (\IV{}) design exploits
prior threshold crossings to identify the effect of past \IA{} receipt on
subsequent NFIP damage patterns, our proxy for pre-disaster mitigation effort.
\emph{Mechanism B}: a contemporaneous fuzzy regression discontinuity (\RD{})
design estimates the effect of current \IA{} declaration on post-disaster
claiming behavior.
[Results pending data acquisition and analysis.]
\end{abstract}

\newpage
\tableofcontents
\newpage

% ============================================================
\section{Introduction}
\label{sec:intro}
% ============================================================

The United States spends tens of billions of dollars annually on federal
disaster assistance through programs such as FEMA Individual Assistance,
which provides direct grants to households affected by presidentially declared
major disasters.
A long-standing concern in the disaster risk management literature is that
public disaster aid, by reducing the expected private cost of unmitigated
exposure, may discourage household investment in mitigation—the canonical
Peltzman~\citeyearpar{peltzman1975} moral hazard mechanism applied to
catastrophe insurance \citep{kunreuther2004,browne2008}.
The empirical magnitude of this effect is an open question with direct policy
relevance: if aid crowds out mitigation, the net welfare effect of expanding
\IA{} coverage could be negative.

The primary empirical challenge is that \IA{} exposure is not randomly assigned:
disaster severity, community characteristics, and political economy all influence
both who receives aid and who later mitigates.
We address this identification problem using the per-capita damage threshold that
FEMA uses to determine \IA{} eligibility under 44 CFR \S{}206.48, which creates
sharp variation in aid receipt as a function of a running variable that is
plausibly exogenous to household mitigation behavior within narrow neighborhoods
of the threshold.

We contribute to a growing quasi-experimental literature on disaster aid and
mitigation behavior
\citep{gallagher2014,horn2021,andor2020,botzen2019,davlasheridze2019}.
Relative to \citeauthor{gallagher2014}'s \citeyearpar{gallagher2014} seminal
work using NFIP purchase decisions as the mitigation proxy, our design:
(i)~instruments prior \IA{} receipt using past threshold crossings rather than
relying on within-ZIP panel variation;
(ii)~uses the contents-to-structural damage ratio—a behavioral outcome that
reflects pre-storm mitigation choices—rather than insurance purchase;
and (iii)~applies a TTR adjustment to the running variable to control for
cross-state variation in state fiscal capacity.

[Remainder of introduction: roadmap, preview of results, outline.]

% ============================================================
\section{Institutional Background}
\label{sec:background}
% ============================================================

\subsection{FEMA Individual Assistance}

FEMA's Individual Assistance program provides direct grants to disaster-affected
households for temporary housing, home repair, and other unmet needs following a
presidential major disaster declaration.
Whether a county receives an \IA{} declaration depends on FEMA's assessment of
whether state resources are ``overwhelmed,'' formalized in a per-capita damage
threshold that has been updated periodically through Federal Register notices
under 44 CFR \S{}206.48.
[Detail on threshold history, CPI adjustments, state vs. county thresholds.]

\subsection{National Flood Insurance Program}

[NFIP structure, mandatory purchase requirement, contents vs. building coverage,
claim payment process, mitigation-relevant behavioral margins.]

\subsection{TTR Adjustment}
\label{sec:ttr}

[Tax-to-revenue ratio adjustment rationale: cross-state fiscal capacity variation
means the same per-capita damage number implies different levels of ``overwhelmed''
for a wealthy vs. a poor state. Three specifications (linear, log, rank) and
selection via first-stage F.]

% ============================================================
\section{Conceptual Framework}
\label{sec:framework}
% ============================================================

\subsection{Mechanism A: Historical Mitigation Response}

Let $m_{i,t}$ denote household $i$'s mitigation investment in period $t$,
and let $A_{i,t-k}$ denote receipt of \IA{} in a prior disaster $k$ periods ago.
The Peltzman hypothesis predicts $\partial m / \partial A < 0$: prior aid receipt
reduces current mitigation.
The complementarity hypothesis \citep{andor2020} predicts $\partial m / \partial A > 0$:
aid enables recovery investment that raises the value of the home and thus the
marginal benefit of mitigation.

\subsection{Mechanism B: Post-Disaster Claiming Response}

Let $c_{i,t}$ denote household $i$'s post-disaster claiming behavior (application
rate, award amount, time to application) following current disaster $t$.
If prior aid creates moral hazard, we expect $\partial c / \partial A_{t} > 0$
for households in counties that just crossed the declaration threshold.

\subsection{Theoretical Predictions}
[Table of sign predictions under Peltzman vs. complementarity, by outcome and mechanism.]

% ============================================================
\section{Data}
\label{sec:data}
% ============================================================

\subsection{FEMA Declarations and Individual Assistance}
[OpenFEMA: DisasterDeclarationsSummaries, HousingAssistanceOwners/Renters,
IndividualAssistanceHousingRegistrantsLargeDisasters. Sample period, geographic coverage.]

\subsection{NFIP Claims}
[FimaNfipClaims: contents and building paid amounts, water depth, date, ZIP.
Sample construction: matched to disaster by disasterNumber then date+county window.]

\subsection{Public Assistance Damage Estimates}
[PublicAssistanceApplicantsPrograms: project costs aggregated to county level
to construct the running variable denominator.]

\subsection{Hazard Intensity Controls}
[USGS STN high-water marks (IDW interpolated to ZIP), HURDAT2 wind radii,
NWS IEM warning lead times.]

\subsection{ACS Demographics}
[ZCTA-level 5-year ACS, 2011--2023 vintages. Key variables: population, income,
owner-occupancy, housing age.]

\subsection{Treasury Tax-to-Revenue Ratios}
[Annual state-level TTR data. Percentile rank within fiscal year used in rank spec.]

\subsection{Sample Construction and Summary Statistics}

\input{../tables/tab_desc_stats_a.tex}
\input{../tables/tab_desc_stats_b.tex}

% ============================================================
\section{Identification Strategy}
\label{sec:identification}
% ============================================================

\subsection{Running Variable Construction}
\label{sec:running_var}

The running variable for the contemporaneous design (Mechanism B) is:
\begin{equation}
  RV_{c,t} = \frac{PA\_damage_{c,t} / pop_{c,t}}{TTR_{s,t} / 100} - \tau_{s,t}
  \label{eq:rv_linear}
\end{equation}
where $PA\_damage_{c,t}$ is county-level Public Assistance project costs,
$pop_{c,t}$ is county population from the nearest prior ACS vintage,
$TTR_{s,t}$ is state $s$'s tax-to-revenue ratio as a percentage of the national
average, and $\tau_{s,t}$ is the applicable per-capita \IA{} threshold for
year $t$.

Two alternative TTR adjustment specifications are also estimated (log and
percentile-rank transformations; see Appendix~\ref{app:ttr}) and the preferred
specification is selected by highest first-stage $F$-statistic
(Table~\ref{tab:first_stage}).

For Mechanism A, the prior running variable $RV_{c,t-k}$ is constructed
analogously using the prior disaster.

\subsection{Mechanism A: Historical IV Design}

The first-stage regression is:
\begin{equation}
  A_{i,t-k} = \alpha_0 + \alpha_1 \mathbf{1}[RV_{c,t-k} \geq 0]
              + g(RV_{c,t-k}) + X_{i,t} + \delta_s + \lambda_t + \varepsilon_{i,t}
  \label{eq:fs_a}
\end{equation}
and the second-stage structural equation is:
\begin{equation}
  Y_{i,t} = \beta_0 + \beta_1 \hat{A}_{i,t-k}
              + g(RV_{c,t-k}) + X_{i,t} + H_{c,t} + \delta_s + \lambda_t + u_{i,t}
  \label{eq:ss_a}
\end{equation}
where $Y_{i,t}$ is a mitigation proxy outcome, $g(\cdot)$ is a linear control
for the running variable, $X_{i,t}$ are ACS demographics, $H_{c,t}$ are hazard
intensity controls, and $\delta_s$, $\lambda_t$ are state and year fixed effects.
Standard errors are clustered at the county level.

\subsection{Mechanism B: Fuzzy RD Design}

The sharp RD estimate of the intent-to-treat effect is:
\begin{equation}
  \tau^{ITT} = \lim_{r \downarrow 0} \mathbb{E}[Y | RV=r] -
               \lim_{r \uparrow 0} \mathbb{E}[Y | RV=r]
\end{equation}
The fuzzy RD (LATE) scales by the first-stage jump in \IA{} declaration
probability:
\begin{equation}
  \tau^{LATE} = \frac{\tau^{ITT}}{\lim_{r \downarrow 0} \Pr[\IA=1|RV=r] -
                                   \lim_{r \uparrow 0} \Pr[\IA=1|RV=r]}
\end{equation}
We estimate both using local linear regression with triangular kernel and
\citet{calonico2014} optimal bandwidth (CCT).

\subsection{Validity Tests}

\subsubsection{McCrary Density Test}
Figure~\ref{fig:mccrary_a} and Table~\ref{tab:mccrary} report \citet{cattaneo2020}
density discontinuity tests at the threshold for both mechanisms.

\subsubsection{Balance on Predetermined Covariates}
Table~\ref{tab:balance} reports local linear \RD{} estimates of the discontinuity
in five predetermined covariates (median household income, presidential vote margin,
owner-occupancy rate, pre-1980 housing share, county population).
We apply a Bonferroni correction for $k=5$ tests.

\subsubsection{Placebo Thresholds}
Table~\ref{tab:placebo_thresholds} reports \RD{} estimates at placebo thresholds
located at the 10th, 25th, 75th, and 90th percentiles of the running variable
distribution, excluding observations near the true threshold.

% ============================================================
\section{Results}
\label{sec:results}
% ============================================================

\subsection{First Stage}
\label{sec:first_stage}

\input{../tables/tab_first_stage.tex}

\begin{figure}[htbp]
  \centering
  \includegraphics[width=0.8\textwidth]{fig_binscatter_a.png}
  \caption{First Stage — Mechanism A: Prior IA receipt vs. prior running variable.}
  \label{fig:binscatter_a}
\end{figure}

\begin{figure}[htbp]
  \centering
  \includegraphics[width=0.8\textwidth]{fig_binscatter_b.png}
  \caption{First Stage — Mechanism B: IA declaration probability vs. running variable.}
  \label{fig:binscatter_b}
\end{figure}

\subsection{Mechanism A: Effect of Prior IA on Pre-Disaster Mitigation}
\label{sec:mech_a}

\input{../tables/tab_mech_a_2sls.tex}

\subsection{Mechanism B: Effect of IA Declaration on Post-Disaster Claiming}
\label{sec:mech_b}

\input{../tables/tab_mech_b_rd.tex}

\subsection{SDID Estimates}
\label{sec:sdid}

\input{../tables/tab_sdid.tex}

% ============================================================
\section{Robustness}
\label{sec:robustness}
% ============================================================

\input{../tables/tab_robustness.tex}

\begin{figure}[htbp]
  \centering
  \includegraphics[width=0.7\textwidth]{fig_donut_rd.png}
  \caption{Donut RD Sensitivity — Mechanism B.}
  \label{fig:donut}
\end{figure}

\begin{figure}[htbp]
  \centering
  \includegraphics[width=0.9\textwidth]{fig_ttr_spec_comparison.png}
  \caption{Coefficient stability across TTR specifications.}
  \label{fig:ttr_spec}
\end{figure}

% ============================================================
\section{Heterogeneity}
\label{sec:heterogeneity}
% ============================================================

\input{../tables/tab_heterogeneity.tex}

\begin{figure}[htbp]
  \centering
  \includegraphics[width=0.75\textwidth]{fig_temporal_decay.png}
  \caption{Mechanism A: Temporal decay of the treatment effect by years since
           prior IA event (bins pre-specified per \citealt{gallagher2014}).}
  \label{fig:decay}
\end{figure}

\begin{figure}[htbp]
  \centering
  \includegraphics[width=0.75\textwidth]{fig_tenure_interaction.png}
  \caption{Mechanism A: Treatment effect by owner-occupancy tercile.}
  \label{fig:tenure}
\end{figure}

% ============================================================
\section{Discussion}
\label{sec:discussion}
% ============================================================

[Interpretation of main results under Peltzman vs. complementarity framing.
Comparison with Gallagher (2014), Andor et al. (2020), Botzen et al. (2019).
Welfare implications for IA program design.
Limitations: running variable fuzziness, NFIP coverage gaps, PA damage measurement.]

% ============================================================
\section{Conclusion}
\label{sec:conclusion}
% ============================================================

[Summary of findings. Policy implications for FEMA disaster aid design.
Directions for future work: heterogeneous effects by income, dynamic panel,
extension to other IA-adjacent programs.]

% ============================================================
%% Appendices
% ============================================================
\appendix

\section{TTR Adjustment Specifications}
\label{app:ttr}

The three TTR-adjusted running variable specifications (pre-specified in the
project plan and selected by first-stage $F$-statistic) are:
\begin{align}
  RV^{\text{linear}}_{c,t} &= \frac{PC_{c,t}}{TTR_{s,t}/100} - \tau_{s,t}
  \label{eq:rv_linear_app} \\
  RV^{\text{log}}_{c,t}    &= PC_{c,t} \cdot \ln\!\left(\frac{100}{TTR_{s,t}}\right) - \tau_{s,t}
  \label{eq:rv_log} \\
  RV^{\text{rank}}_{c,t}   &= PC_{c,t} \cdot \bigl(1 + (1 - \text{rank}_{s,t})\bigr) - \tau_{s,t}
  \label{eq:rv_rank}
\end{align}
where $PC_{c,t}$ is per-capita PA project costs, $TTR_{s,t}$ is the state
tax-to-revenue ratio as a percentage of the national average, $\text{rank}_{s,t}$
is its within-year percentile rank, and $\tau_{s,t}$ is the applicable
per-capita threshold.

\section{McCrary Density Test Results}
\label{app:mccrary}

\input{../tables/tab_mccrary.tex}

\begin{figure}[htbp]
  \centering
  \begin{subfigure}[t]{0.48\textwidth}
    \includegraphics[width=\textwidth]{fig_mccrary_a.png}
    \caption{Mechanism A (prior running variable)}
    \label{fig:mccrary_a}
  \end{subfigure}
  \hfill
  \begin{subfigure}[t]{0.48\textwidth}
    \includegraphics[width=\textwidth]{fig_mccrary_b.png}
    \caption{Mechanism B (contemporaneous running variable)}
    \label{fig:mccrary_b}
  \end{subfigure}
  \caption{McCrary density test plots.
           Bars: binned density estimate.
           Lines: local linear fit on each side.
           Null hypothesis: density is continuous at threshold.}
\end{figure}

\section{Balance Tests}
\label{app:balance}

\input{../tables/tab_balance.tex}

\begin{figure}[htbp]
  \centering
  \includegraphics[width=0.75\textwidth]{fig_balance_coefplot.png}
  \caption{RD balance coefficient plot. Red markers indicate significance
           at Bonferroni-adjusted $\alpha = 0.01$.}
  \label{fig:balance}
\end{figure}

\section{Anderson-Rubin Confidence Intervals}
\label{app:ar}

If the first-stage $F$-statistic falls below 10, we report weak-instrument-robust
Anderson-Rubin \citeyearpar{andersonrubin1949} confidence intervals in addition
to the standard 2SLS estimates.

\IfFileExists{../tables/tab_mech_a_ar_ci.tex}{%
  \input{../tables/tab_mech_a_ar_ci.tex}%
}{%
  [AR CI table will appear here if first-stage F\,$<$\,10 in the data.]%
}

\section{Gate Check Report}
\label{app:gate}

The gate check report produced by \texttt{analysis/00\_gate\_check.py} is
reproduced below for transparency. Analysis proceeds only if all five gates pass.

\begin{verbatim}
[Gate check report — paste contents of data/analysis/gate_check_report.txt here]
\end{verbatim}

% ============================================================
%% References
% ============================================================
\bibliographystyle{aer}
\bibliography{references}

\end{document}
